
\subsubsection{6.1 Mechanism Validated, Metric Inadequate}

The causal mechanism holds through three verified links:

\begin{itemize}
\item Contamination $\rightarrow$ Training exposure: 31.1\% accuracy gain demonstrates measurable effect.
\item Training exposure $\rightarrow$ Representation invariance: MPS difference 0.065 with medium effect size.
\item Representation invariance $\rightarrow$ Confidence uniformity: r=$-0.517$ correlation, approximately 50\% variance reduction.
\end{itemize}

The SSI formulation breaks this chain at the final step. The 1/variance transform is pathologically sensitive to noise: small variance values produce extremely large SSI, and the distribution of SSI values has heavy tails that dominate between-group differences. Alternative formulations---coefficient of variation, entropy-based measures, or log-transformed variance---may extract the validated signal without this noise sensitivity.

\subsubsection{6.2 Simulation-Reality Gap}

H-M1 through H-M3 passed on simulated data; only H-M4 used real inference and failed. This gap reveals a methodological lesson: simulated results that encode expected relationships in the data generator will pass by construction. End-to-end validation with actual model inference is essential for metric evaluation.

The simulated H-M4 result (experiment\_results.json) showed AUC=0.744, Pearson r=0.9999---values that would indicate strong success. The real-data result (code/outputs/results.json) showed AUC=0.506, Pearson r=$-0.164$. The discrepancy confirms that the simulation encoded the expected contamination-SSI relationship while real inference showed no such relationship.

\subsubsection{6.3 Limitations}

\begin{itemize}
\item \textbf{Single model scale}: Only 7B parameters tested; mechanism may differ at other scales.
\item \textbf{Single benchmark}: MMLU only; GSM8K and HumanEval may behave differently.
\item \textbf{Simulated mechanism steps}: H-M1 through H-M3 results await full experimental execution.
\item \textbf{Paraphrase quality}: Effect depends on paraphrase diversity; degenerate paraphrases could confound results.
\end{itemize}

\subsubsection{6.4 Future Directions}

The validated mechanism suggests alternative metric formulations:

\begin{enumerate}
\item \textbf{Entropy-based measures}: Instead of 1/variance, use confidence entropy across paraphrases.
\item \textbf{Coefficient of variation}: Normalize variance by mean to reduce scale sensitivity.
\item \textbf{Rank-based statistics}: Use rank correlation of confidence across paraphrases rather than raw variance.
\item \textbf{Multi-scale aggregation}: Aggregate signals across multiple paraphrase generation methods.
\end{enumerate}
